%
% TU/e Style Master Thesis template for LaTeX
%
% 2021 Crated by Marko Boon to new corporate identity,
% Based on a template by Thijs Nugteren and Joos Buijs.
%
% THIS IS THE MAIN FILE (i.e. compile this file, compiling the others directly won't work)
%
\documentclass[a4paper,11pt,twoside]{book}
% \documentclass[a4paper]{article}
 

\usepackage[backref=true,style=alphabetic]{biblatex} 
\addbibresource{references.bib}
\usepackage{tuethesis2021}




\usepackage[
    contentBlocks,
    debugExtensions,
    definitionLists,
    fancy_lists,
    fencedCode,
    hashEnumerators,
    inlineNotes,
    jekyllData,
    notes,
    pipeTables,
    rawAttribute,
    headerAttributes,
    smartEllipses,
    strikeThrough,
    subscripts,
    superscripts,
    tableCaptions,
    taskLists,
    citations,
    hybrid,
    relativeReferences,
    texComments,
]{markdown} 

\usepackage{nomencl} 
% \usepackage{kpathsea}

% \begin{filecontents}[overwrite,noheader]{markdown-languages.json}
%     {
%       "nix": "JavaScript Object Notation"
%     }
% \end{filecontents}
    

%
% These commands need to be defined in order to produce a correct and personalized document.
% You will get error messages if not all these commands are defined. 
% If you don't need a specific command, just leave the argument empty. 
% For example: \docsubtitle{}
%     
\newcommand{\doctitle}{Secure Dynamic Setup for MPyC Sessions to Support (Ad Hoc) Multiparty Computation}
\newcommand{\docsubtitle}{Preparation phase report}
\newcommand{\me}{Emil Nikolov}
\newcommand{\studentNumber}{0972305}
\newcommand{\email}{emil.e.nikolov@gmail.com}
\newcommand{\version}{Draft version}
\newcommand{\placeMonthYear}{Eindhoven, November 2022}
\newcommand{\department}{Department of Mathematics and Computer Science}
\newcommand{\group}{Coding Theory and Cryptology Group}
\newcommand{\firstCommitteeMember}{Dr. ir. L.A.M. (Berry) Schoenmakers} % use all the titles for your committee members!
\newcommand{\secondCommitteeMember}{Dr. Mike Holenderski} % usually the daily supervisor
\newcommand{\thirdCommitteeMember}{Dr. Savio Sciancalepore} % usually the external member
\newcommand{\builddate}{\today} 
\def\markdownOptionOutputDir{build}
\makenomenclature 
\renewcommand{\nomname}{List of Abbreviations}

\usepackage{mdframed}
% \usepackage{pygmentize}
\usepackage{minted}
% \usemintedstyle{xcode}
% \usemintedstyle{paraiso-light}
\usemintedstyle{arduino}

\usepackage{csquotes}
\usepackage[english]{babel} 

\begin{document}




\definecolor{LightGray}{gray}{0.9}

% % \surroundwithmdframed[bgcolor=DarkGray]{minted}
% \BeforeBeginEnvironment{minted}{\begin{mdframed} bgcolor=DarkGray}
%         \AfterEndEnvironment{minted}{\end{mdframed}}

% \begin{minted}[bgcolor=darkgray]{python}
% def test :
%     print 123

% \end{minted}


% \begin{titlepage}
    \begin{center}
        \includegraphics[height=2cm]{figures/TUe-logo-descriptor-line-scarlet-rgb}\\
        %\LARGE
        %Eindhoven University of Technology \\
        \large
        \mbox{}\department\\\mbox{}\group

        % \vspace*{10cm}
        \vspace{1.6cm}

        \setlength{\TPHorizModule}{1mm}
        \setlength{\TPVertModule}{\TPHorizModule}
        % Set the Paragraph Indent to zero, so the first line is not Indented
        % Back-up the current value so it can be put back at the end of the title page
        \newlength{\backupparindent}
        \setlength{\backupparindent}{\parindent}
        \setlength{\parindent}{0mm}
        % Begins a textbox at 72 mm from the left of the edge of the paper and 89 mm from the top
        % The width of the textbox is 95 mm (167 - 72 mm)
        % The height of the box cannot be defined, so it is your task to keep the text not too long
        % \begin{textblock}{95}(62,89)
        \vspace*{1mm}
        \huge
        \textbf{\doctitle \\}
        \Large
        \vspace*{5mm}
        \textit{\docsubtitle}\\
        \vspace*{10mm}
        \Large
        \me\\
        % \end{textblock}

        \large
        Supervisors:\\
        % \begin{tabular}{r}
        \mbox{}\firstCommitteeMember  \\
        \mbox{}\secondCommitteeMember \\
        \mbox{}\thirdCommitteeMember  \\
        % \end{tabular}

        \vfill
        \version

        \vfill
        %\docdate \\
        \large
        \placeMonthYear\\

        % Put the Paragraph Indent back to its original value
        \setlength{\parindent}{\backupparindent}
    \end{center}
\end{titlepage}
\normalsize


\begin{titlepage}
    \begin{center}
        \vspace{1.6cm}
        \includegraphics[height=2cm]{figures/TUe-logo-descriptor-line-scarlet-rgb}\\
        \vspace{1.6cm}

        \Large
        \department\\
        \group\\
        \vspace*{1cm}
        \Huge
        \textbf{\doctitle}
        \vspace{0.5cm}

        \LARGE
        \docsubtitle
        \vspace{1.5cm}

        \large
        \textbf{\me}\\
        \vspace{0.4cm}
        Id nr: \studentNumber \\
        \texttt{\email}
        \vfill


        Supervisor : \firstCommitteeMember\\

        \vfill

        Committee members:\\
        \firstCommitteeMember\\
        \secondCommitteeMember\\
        \thirdCommitteeMember\\
        \vfill
        \builddate \\

    \end{center}
\end{titlepage}

% The Flexible Spin-Lock Model (FSLM) is a real-time multiprocessor resource sharing protocol.
% Earlier work has been carried out in  defining the protocol, preliminary schedulability analysis and an initial implementation of the FSLM
% \nomenclature{FSLM}{Flexible Spin-Lock Model}.
% The aim of this assignment is to develop schedulability analysis for FSLM taking into consideration the implementation overheads in the Erika  enterprise RTOS
% \nomenclature{RTOS}{Real Time Operating System}
% instantiated on an Altera Nios II platform with 4 soft cores. The assignment also aims at exploring the effects of improvements on the existing implementation. The purpose of this report is to detail the preparatory steps carried out to perform the assignment.


% The title page will be inserted automatically, here.

\clearpage

%Sometimes line numbers are nice, uncomment the next line to enable:
%\linenumbers


% \chapter*{Preface}\label{chapter:preface}

Please write all your preface text here. If you do so, don't forget to thank your supervisor, other committee members, your family, colleagues etc.\ etc. 


% \printnomenclature[5em]


% \listoffigures

% \listoftables

% \lstlistoflistings

\mainmatter

% \chapter*{Abstract}\label{chapter:abstract}




% \begin{abstract}
% \end{abstract}

\chapter*{Abstract}
\markdownInput{notion/00-abstract.md}
% \markdownInput[slice=^ ^section2]{notion/00-abstract.md}


% \begin{markdown}
% \in
% \end{markdown}
% {notion/00-abstract.md}

\tableofcontents


\markdownInput{notion/01-introduction.md}

\markdownInput{notion/02-evaluation-setup.md}
\markdownInput{notion/03-evaluation-setup-implementation.md}
\markdownInput{notion/04-planning.md}
\markdownInput{notion/05-conclusions.md}

% \chapter{Project goals}\label{chapter:first_real_chapter}

\section{List of goals of the assignment}
\subsection{Implementation Goals (I.G)}
\subsection{Measurement Goals (M.G)}
\subsection{Analysis Goals (A.G)}
\subsection{Documentation Goals (D.G)}
\section{Identified risks}
\section{Scope and limitations}
\section{Research questions}


% \chapter{Literature survey}\label{chapter:second_real_chapter}



% \chapter{Conclusions}\label{chapter:conclusions}

Write your conclusions here. 



% \bibliographystyle{plain}
\printbibliography
% \bibliographystyle{plainurl}
% \bibliography{references}

%\appendix
%\addcontentsline{toc}{chapter}{Appendix}


% \chapter{My First Appendix}
In this file (appendices/main.tex) you can add appendix chapters, just as you did in the thesis.tex file for the `normal' chapters.
You can also choose to include everything in this single file, whatever you prefer.

\end{document}
